
\subsubsection{Peripheral}

The Peripheral MPU manages the 39 peripheral modules. This MMU can be configured per peripheral to only allow access from a process with a certain PID. The registers to configure this are detailed in Table \ref{table:MPU_for_Peripheral}.

\begin{longtable}[c]{|p{4cm}|p{2cm}|p{9cm}|}
    
    \caption{MPU for Peripheral}
    \label{table:MPU_for_Peripheral}
    \\\hline
    \rowcolor{lightgray}
    & \multicolumn{2}{c|}{Authority} \\
    \cline{2-3}
    \rowcolor{lightgray}
    \multirow{-2}*{Peripheral} & PID = 0/1 & PID = 2 \verb+~+ 7 \\
    \hline
    \endfirsthead
    
    \hline
    \rowcolor{lightgray}
    & \multicolumn{2}{c|}{Authority} \\
    \cline{2-3}
    \rowcolor{lightgray}
    \multirow{-2}*{Peripheral} & PID = 0/1 & PID = 2 \verb+~+ 7 \\
    \hline
    \endhead

    \endfoot
    
    \hline
    \endlastfoot
    
    DPort Register        & Access & Forbidden \\
    \hline
    AES Accelerator       & Access & Forbidden \\
    \hline
    RSA Accelerator       & Access & Forbidden \\
    \hline
    SHA Accelerator       & Access & Forbidden \\
    \hline
    Secure Boot           & Access & Forbidden \\
    \hline
    
    
    Cache MMU Table       & Access & Forbidden \\
    \hline
    PID Controller        & Access & Forbidden \\
    \hline
    UART0                 & Access & DPORT\_AHBLITE\_MPU\_TABLE\_UART\_REG \\
    \hline
    SPI1                  & Access & DPORT\_AHBLITE\_MPU\_TABLE\_SPI1\_REG \\
    \hline
    SPI0                  & Access & DPORT\_AHBLITE\_MPU\_TABLE\_SPI0\_REG \\
    \hline
    GPIO                  & Access & DPORT\_AHBLITE\_MPU\_TABLE\_GPIO\_REG \\
    \hline
    
    
    RTC                   & Access & DPORT\_AHBLITE\_MPU\_TABLE\_RTC\_REG \\
    \hline
    IO MUX                & Access & DPORT\_AHBLITE\_MPU\_TABLE\_IO\_MUX\_REG \\
    \hline
    
    
    SDIO Slave            & Access & DPORT\_AHBLITE\_MPU\_TABLE\_HINF\_REG \\
    \hline
    UDMA1                 & Access & DPORT\_AHBLITE\_MPU\_TABLE\_UHCI1\_REG \\
    \hline
    
    
    I2S0                  & Access & DPORT\_AHBLITE\_MPU\_TABLE\_I2S0\_REG \\
    \hline
    UART1                 & Access & DPORT\_AHBLITE\_MPU\_TABLE\_UART1\_REG \\
    \hline
    
    
    I2C0                  & Access & DPORT\_AHBLITE\_MPU\_TABLE\_I2C\_EXT0\_REG \\
    \hline
    UDMA0                 & Access & DPORT\_AHBLITE\_MPU\_TABLE\_UHCI0\_REG \\
    \hline
    SDIO Slave            & Access & DPORT\_AHBLITE\_MPU\_TABLE\_SLCHOST\_REG \\
    \hline
    RMT                   & Access & DPORT\_AHBLITE\_MPU\_TABLE\_RMT\_REG \\
    \hline
    PCNT                  & Access & DPORT\_AHBLITE\_MPU\_TABLE\_PCNT\_REG \\
    \hline
    SDIO Slave            & Access & DPORT\_AHBLITE\_MPU\_TABLE\_SLC\_REG \\
    \hline
    LED PWM               & Access & DPORT\_AHBLITE\_MPU\_TABLE\_LEDC\_REG \\
    \hline
    Efuse Controller      & Access & DPORT\_AHBLITE\_MPU\_TABLE\_EFUSE\_REG \\
    \hline
    Flash Encryption      & Access & DPORT\_AHBLITE\_MPU\_TABLE\_SPI\_ENCRYPT\_REG \\
    \hline
    
    
    PWM0                  & Access & DPORT\_AHBLITE\_MPU\_TABLE\_PWM0\_REG \\
    \hline
    TIMG0                 & Access & DPORT\_AHBLITE\_MPU\_TABLE\_TIMERGROUP\_REG \\
    \hline
    TIMG1                 & Access & DPORT\_AHBLITE\_MPU\_TABLE\_TIMERGROUP1\_REG \\
    \hline
    SPI2                  & Access & DPORT\_AHBLITE\_MPU\_TABLE\_SPI2\_REG \\
    \hline
    SPI3                  & Access & DPORT\_AHBLITE\_MPU\_TABLE\_SPI3\_REG \\
    \hline
    SYSCON                & Access & DPORT\_AHBLITE\_MPU\_TABLE\_APB\_CTRL\_REG \\
    \hline
    I2C1                  & Access & DPORT\_AHBLITE\_MPU\_TABLE\_I2C\_EXT1\_REG \\
    \hline
    SDMMC                 & Access & DPORT\_AHBLITE\_MPU\_TABLE\_SDIO\_HOST\_REG \\
    \hline
    EMAC                  & Access & DPORT\_AHBLITE\_MPU\_TABLE\_EMAC\_REG \\
    \hline
    
    
    PWM1                  & Access & DPORT\_AHBLITE\_MPU\_TABLE\_PWM1\_REG \\
    \hline
    I2S1                  & Access & DPORT\_AHBLITE\_MPU\_TABLE\_I2S1\_REG \\
    \hline
    UART2                 & Access & DPORT\_AHBLITE\_MPU\_TABLE\_UART2\_REG \\
    \hline
    
    
    RNG               & Access & DPORT\_AHBLITE\_MPU\_TABLE\_PWR\_REG \\
    
\end{longtable}


Each bit of register DPORT\_AHBLITE\_MPU\_TABLE\_\textcolor{red}{\it{X}}\_REG determines whether each process can access the peripherals managed by the register. For details please see Table \ref{table:DPORT_AHBLITE_MPU_TABLE_X_REG}.
When a bit of register DPORT\_AHBLITE\_\\MPU\_TABLE\_\textcolor{red}{\it{X}}\_REG is 1, it means that a process with the corresponding PID can access the corresponding peripheral of the register. Otherwise, the process cannot access the corresponding peripheral.



\begin{table}[H]
    \centering
    \caption{DPORT\_AHBLITE\_MPU\_TABLE\_\textcolor{red}{\it{X}}\_REG}
    \label{table:DPORT_AHBLITE_MPU_TABLE_X_REG}
       \begin{tabular}{|p{11cm}|p{4cm}|}
        \hline
        \rowcolor{lightgray}
        PID & 2 3 4 5 6 7 \\
        \hline
        DPORT\_AHBLITE\_MPU\_TABLE\_\textcolor{red}{\it{X}}\_REG bit & 0 1 2 3 4 5 \\
        \hline
    \end{tabular}
\end{table}

All the DPORT\_AHBLITE\_MPU\_TABLE\_\textcolor{red}{\it{X}}\_REG registers are in peripheral DPort Register. Only processes with PID 0/1 can modify these registers.